\medskip\noindent
Let $\mathcal{A}$ be a Grothendieck abelian category and let $U$ be a
generator for $\mathcal{A}$, see
Definition \ref{definition-grothendieck-conditions}.
Let $R = \Hom_\mathcal{A}(U, U)$. Consider the functor
$G : \mathcal{A} \to \text{Mod}_R$ given by
$$
G(A) = \Hom_\mathcal{A}(U, A)
$$
endowed with its canonical right $R$-module structure.

\begin{lemma}
\label{lemma-gabriel-popescu-left-adjoint}
The functor $G$ above has a left adjoint
$F : \text{Mod}_R \to \mathcal{A}$.
\end{lemma}

\begin{proof}
We will give two proofs of this lemma.

\medskip\noindent
The first proof will use the adjoint functor theorem, see
Categories, Theorem \ref{categories-theorem-adjoint-functor}.
Observe that $G : \mathcal{A} \to \text{Mod}_R$ is left exact and sends
products to products. Hence $G$ commutes with limits. To check the set
theoretical condition in the theorem, suppose that $M$ is an object of
$\text{Mod}_R$. Choose a suitably large cardinal $\kappa$ and denote $E$
a set of objects of $\mathcal{A}$ such that every object $A$ with
$|A| \leq \kappa$ is isomorphic to an element of $E$. This is possible
by Lemma \ref{lemma-set-iso-classes-bounded-size}. Set
$I = \coprod_{A \in E} \Hom_R(M, G(A))$.
We think of an element $i \in I$ as a pair $(A_i, f_i)$.
Finally, let $A$ be an arbitrary object of $\mathcal{A}$
and $f : M \to G(A)$ arbitrary. We are going to think of
elements of $\Im(f) \subset G(A) = \Hom_\mathcal{A}(U, A)$
as maps $u : U \to A$. Set
$$
A' = \Im(\bigoplus\nolimits_{u \in \Im(f)} U \xrightarrow{u} A)
$$
Since $G$ is left exact, we see that $G(A') \subset G(A)$
contains $\Im(f)$ and we get $f' : M \to G(A')$ factoring $f$.
On the other hand, the object $A'$ is
the quotient of a direct sum of at most $|M|$ copies of $U$.
Hence if $\kappa = |\bigoplus_{|M|} U|$, then we see that $(A', f')$
is isomorphic to an element $(A_i, f_i)$ of $E$ and we conclude that $f$
factors as $M \xrightarrow{f_i} G(A_i) \to G(A)$ as desired.

\medskip\noindent
The second proof will give a construction of $F$ which will show
that ``$F(M) = M \otimes_R U$'' in some sense. Namely, for any
$R$-module $M$ we can choose a resolution
$$
\bigoplus\nolimits_{j \in J} R \to
\bigoplus\nolimits_{i \in I} R \to
M \to 0
$$
Then we define $F(M)$ by the corresponding exact sequence
$$
\bigoplus\nolimits_{j \in J} U \to
\bigoplus\nolimits_{i \in I} U \to
F(M) \to 0
$$
This construction is independent of the choice of the resolution
and is functorial; we omit the details.
For any $A$ in $\mathcal{A}$ we obtain an exact sequence
$$
0 \to \Hom_\mathcal{A}(F(M), A) \to
\prod\nolimits_{i \in I} G(A) \to
\prod\nolimits_{j \in J} G(A)
$$
which is isomorphic to the sequence
$$
0 \to \Hom_R(M, G(A)) \to
\Hom_R(\bigoplus\nolimits_{i \in I} R, G(A)) \to
\Hom_R(\bigoplus\nolimits_{j \in J} R, G(A))
$$
which shows that $F$ is the left adjoint to $G$.
\end{proof}
